\section{Experiment Setup}
\label{sec:experiment_setup}

Section~\ref{sec:framework} defined the framework abstractly. This
section instantiates it experimentally: which SLMs sit at which stages
under the pre-registered DOE (Section~\ref{sec:experiment_setup_doe}),
which benchmark functions the framework is evaluated on
(Section~\ref{sec:experiment_setup_datasets}), how statistical inference
is drawn (Section~\ref{sec:experiment_setup_stats}), and how the
strict-AND policy is tested for robustness under threshold perturbation
(Section~\ref{sec:experiment_setup_sensitivity}). We also describe the
post-processing categorization of every closure row into decision
reasons (Section~\ref{sec:experiment_setup_postproc}) and the capability
preservation protocol used to compare hetero and mono cells at
per-sample resolution (Section~\ref{sec:experiment_setup_capability}).
The human-evaluation protocol used to validate the framework's judge
signal against human raters is specified in
Section~\ref{sec:experiment_setup_human_eval}.

\subsection{Design of experiments}
\label{sec:experiment_setup_doe}

A cell is a stage-to-model assignment $a : \mathcal{S} \to \mathcal{M} \cup \{\bot\}$, where $\bot$ denotes an ablated stage. The DOE consists of 20 pre-registered cells across three strata:

\begin{equation}
  \mathcal{C}_{\textsf{DOE}} \;=\;
  \underbrace{\bigl\{a : a(\textsf{spec}) = a(\textsf{test}) = a(\textsf{code})\bigr\}}_{\text{Mono stratum (6 cells)}}
  \;\cup\;
  \underbrace{\bigl\{a : |\{a(s) : s \in \mathcal{S}\}| \geq 2\bigr\}}_{\text{Hetero stratum (11 cells)}}
  \;\cup\;
  \underbrace{\mathcal{C}_{\textsf{null}}}_{\text{Null stratum (3 cells)}}.
  \label{eq:doe}
\end{equation}

\begin{table}[t]
\centering
\small
\caption{Pre-registered 20-cell DOE over six open-weight SLMs and three pipeline stages ($6^3 = 216$ full design space, sampled purposively, not by fractional factorial). Three strata: 6 mono, 11 hetero, 3 null. Hypothesis is the rationale recorded at design time.}
\label{tab:doe_summary}
\begin{tabular}{@{}llllll@{}}
\toprule
Cell & Stratum & $L_{\textsf{spec}}$ & $L_{\textsf{test}}$ & $L_{\textsf{code}}$ & Hypothesis \\
\midrule
M1 & mono & llama3.2 & llama3.2 & llama3.2 & Small-dense self-consistency floor (3\,B Meta). \\
M2 & mono & phi4 & phi4 & phi4 & Mid-dense reasoning-tuned mono baseline (14\,B Microsoft). \\
M3 & mono & qwen3.6 & qwen3.6 & qwen3.6 & Dense mono baseline (27\,B Alibaba). \\
M4 & mono & gemma4 & gemma4 & gemma4 & MoE mono baseline (26\,B Google). \\
M5 & mono & mistral-s.\ 3.2 & mistral-s.\ 3.2 & mistral-s.\ 3.2 & Function-calling-tuned dense mono (24\,B Mistral). \\
M6 & mono & qwen3-coder & qwen3-coder & qwen3-coder & Code-specialised MoE self-consistency (30\,B Alibaba). \\
H1 & hetero & phi4 & qwen3-coder & qwen3-coder & Predicate-reasoner at spec plus coder at test/code. \\
H2 & hetero & qwen3.6 & phi4 & qwen3-coder & Dense spec, reasoner tests, coder synth. \\
H3 & hetero & gemma4 & mistral-s.\ 3.2 & qwen3-coder & Cross-family triple: Google $\to$ Mistral $\to$ Alibaba-coder. \\
H4 & hetero & qwen3-coder & qwen3-coder & llama3.2 & Cheap drafter at synth: does a 3\,B model suffice? \\
H5 & hetero & llama3.2 & qwen3-coder & qwen3-coder & Cheap drafter at spec: can a 3\,B docstring survive downstream? \\
H6 & hetero & phi4 & qwen3.6 & phi4 & Phi sandwich with a different-family middle stage. \\
H7 & hetero & qwen3-coder & qwen3.6 & phi4 & Reverse-capability gradient: strongest first, weakest last. \\
H8 & hetero & gemma4 & qwen3-coder & mistral-s.\ 3.2 & MoE $\to$ MoE $\to$ dense: does architecture matching matter? \\
H9 & hetero & qwen3-coder & qwen3.6 & qwen3-coder & Strong-sandwich: MoE bookends dense for synthesis stability. \\
H10 & hetero & mistral-s.\ 3.2 & phi4 & qwen3-coder & Best-of-each-family-stage: Mistral spec + Phi test + Qwen-coder synth. \\
H11 & hetero & phi4 & gemma4 & qwen3-coder & Phi spec + Gemma test + Qwen-coder synth (alt family triple). \\
N1 & null & llama3.2 & llama3.2 & llama3.2 & Word-shuffled input: expected metric $\approx 0$ if metric is sound. \\
N2 & null & (skip) & qwen3-coder & qwen3-coder & Spec-stage ablation: quantifies $L_{\textsf{spec}}$ contribution. \\
N3 & null & qwen3-coder & (skip) & qwen3-coder & Test-stage ablation: quantifies $L_{\textsf{test}}$ contribution. \\
\bottomrule
\end{tabular}
\end{table}


\begin{table}[t]
\centering
\caption{Small Language Models in the Chapter-3 pipeline. All models are open-weight and served via Ollama; total parameter count is $<30$ B per the §1.1 SLM definition.}
\label{tab:model_lineup}
\begin{tabular}{@{}lccccc@{}}
\toprule
Role & Ollama tag & Family & Size (B) & Architecture & Year \\
\midrule
Pipeline & llama3.2:3b & Meta & 3.0 & dense & 2024 \\
Pipeline & phi4:14b & Microsoft & 14.0 & dense & 2024 \\
Pipeline & qwen3.6:27b & Alibaba & 27.0 & dense & 2026 \\
Pipeline & gemma4:26b & Google & 26.0 & MoE & 2026 \\
Pipeline & mistral-small3.2:24b & Mistral & 24.0 & dense & 2025 \\
Pipeline & qwen3-coder:30b & Alibaba-coder & 30.0 & MoE & 2025 \\
Judge & deepseek-r1:14b & DeepSeek & 14.0 & dense & 2025 \\
\bottomrule
\end{tabular}
\end{table}


Table~\ref{tab:doe_summary} presents the complete DOE with the pre-registered hypothesis for each cell. The full 6-model $\times$ 3-stage design space contains $6^3 = 216$ cells; we sample 20 as a hypothesis-driven purposive selection (not a formal fractional factorial). Hetero cells are permitted two distinct models (H1, H4, H5, H6, H9 use a same-model repetition at two of three stages) or three distinct models (H2, H3, H7, H8, H10, H11), reflecting the design goal of probing specific stage-pairings rather than orthogonalising all six SLMs across all three stages. The mono stratum establishes each SLM's self-consistency floor as a baseline. The hetero stratum assigns SLMs across stages by design-time hypotheses about how model capability might align with per-stage demands (see the Hypothesis column of Table~\ref{tab:doe_summary}); the empirical validation of these hypotheses is the subject of Section~\ref{sec:results}, and the per-operator dominance patterns reported there are findings rather than pre-registered predictions. The null stratum comprises three control cells: N1 (word-shuffled input; expected metric $\approx 0$ if the closure metric is sound), N2 (spec-stage ablated; $L_{\textsf{spec}} = \bot$), and N3 (test-stage ablated; $L_{\textsf{test}} = \bot$).

The six pipeline SLMs (Table~\ref{tab:model_lineup}) are all open-weight and under 30~B parameters, ensuring the entire experiment runs on a single Colab~Pro+ A100/T4 subscription. Five distinct model families are represented, in ascending size order: Meta with Llama-3.2-3B~\citep{Llama32}, Microsoft with Phi-4-14B~\citep{Phi4}, Mistral with Mistral-Small-3.2-24B~\citep{MistralSmall}, Google with Gemma-4-26B~\citep{Gemma4}, and Alibaba with Qwen3.6-27B~\citep{Qwen36} and Qwen3-Coder-30B~\citep{QwenCoder}. DeepSeek-R1:14B~\citep{DeepSeekR1} serves as the family-disjoint judge $L_J$.

\subsection{Datasets}
\label{sec:experiment_setup_datasets}

The primary dataset (\emph{core}) is a stratified 150-function sample from HumanEval (45~functions) and MBPP (105~functions), drawn with seed~42.

Two held-out subsets were prepared for contamination sensitivity analysis: a 25-problem LiveCodeBench~\citep{Jain2025LiveCodeBench} subset with publication date $\geq$ 2024-12-01 (post-training-cutoff for all evaluated SLMs) and a 25-problem HumanEval-Mutated (HEM) subset obtained by applying function-rename, docstring-paraphrase, and parameter-permutation transformations to a stratified draw from HumanEval. Related benchmark alternatives include BigCodeBench~\citep{Zhuo2025BigCodeBench} and CRUXEval~\citep{Gu2025CruxEval}, both released in 2025 and offering complementary decontamination properties. The HEM subset was swept on H1 and M3 (results in Section~\ref{sec:results_contamination}); the LiveCodeBench subset was prepared but blocked at load time by the removal of script-based dataset loaders from the current HuggingFace \texttt{datasets} library (see Limitations, Section~\ref{sec:limitations_external}).

\subsection{Statistical methodology}
\label{sec:experiment_setup_stats}

\subsubsection{Analysis of variance}

For each closure metric $Y$ (kill rate, pass rate, BERTScore stacked across paths), we fit a Type-III analysis-of-variance model

\begin{equation}
  Y_{ij} \;=\; \mu \;+\; \alpha_i \;+\; \beta_j \;+\; \varepsilon_{ij},
  \quad \varepsilon_{ij} \sim \mathcal{N}(0, \sigma^2),
  \label{eq:anova}
\end{equation}

where $\alpha_i$ is the fixed effect of cell $i \in \mathcal{C}_{\textsf{DOE}}$, $\beta_j$ is the fixed effect of sample $j$, and $\varepsilon_{ij}$ is residual error. To further test whether the three closure paths carry non-redundant information about cell performance, we extend Eq.~\ref{eq:anova} to include a path factor and its interaction with cell:

\begin{equation}
  Y_{ijk} \;=\; \mu \;+\; \alpha_i \;+\; \eta_k \;+\; (\alpha \cdot \eta)_{ik} \;+\; \beta_j \;+\; \varepsilon_{ijk},
  \label{eq:anova_interact}
\end{equation}

where $\eta_k$ is the fixed effect of path $k \in \{1, 2, 3\}$ and $(\alpha \cdot \eta)_{ik}$ is the interaction. A significant interaction validates that closure paths are not merely substitutable measurements of the same underlying quantity. We fit Eqs.~\ref{eq:anova}--\ref{eq:anova_interact} via \texttt{statsmodels.formula.api.ols} with Type-III sums of squares (\texttt{anova\_lm(typ=3)}), which tests each effect against a reduced model dropping only that effect --- appropriate for the unbalanced 150-vs-75 mono/hetero design where Type-I (sequential) sums of squares would be order-dependent. Model residual-degree-of-freedom accounting and standard-error propagation are handled by \texttt{statsmodels}. The regeneration script is \texttt{scripts/build\_review\_addenda.py} in the replication package.

\subsubsection{Pairwise cell comparisons}

Pairwise mean differences between cells $i$ and $k$ are tested via Tukey's honest significant difference:

\begin{equation}
  q_{ik}
  \;=\;
  \frac{\bar{Y}_i - \bar{Y}_k}{\sqrt{MS_{\textsf{within}} / n_h}},
  \label{eq:tukey}
\end{equation}

where $n_h$ is the harmonic mean of per-cell sample sizes (mono cells have $n = 150$; hetero and null cells have $n = 75$). Tukey's HSD controls family-wise error at $\alpha = 0.05$ over the $\binom{20}{2} = 190$ pairwise comparisons via the studentized-range distribution of $q_{ik}$; no additional Bonferroni correction is applied.

\subsubsection{Judge--metric correlation}

The Pearson correlation between the automated metric $Y$ and the judge rating $J$ per path,

\begin{equation}
  r_p \;=\; \frac{\textstyle\sum_{r} (Y_r - \bar{Y})(J_r - \bar{J})}{\sqrt{\textstyle\sum_r (Y_r - \bar{Y})^2 \cdot \textstyle\sum_r (J_r - \bar{J})^2}},
  \label{eq:corr}
\end{equation}

quantifies the extent to which each path's automated metric agrees with the judge SLM's semantic-equivalence rating. Rows with parse failures ($J = -1$) are excluded.

\subsubsection{Mixed-effects logit for per-stage decomposition}

To isolate which stage assignment drives closure, we treat each cell as the product of three stage choices and fit a mixed-effects model. Let $V_{ijp} \in \{0, 1\}$ denote whether sample $j$ achieved strict-AND valid closure on path $p$ under cell $i$ (Eq.~\ref{eq:valid_closure_frac} at row resolution). With $\mathbf{d}^{s}_i \in \{0, 1\}^{|\mathcal{M}| - 1}$ denoting the reference-coded dummy vector for the model assigned to stage $s \in \{\textsf{spec}, \textsf{test}, \textsf{code}\}$ under cell $i$ (reference: qwen3.6), and $\mathbf{d}^{\textsf{path}}_p \in \{0, 1\}^{2}$ the analogous path-dummy vector (reference: Path~1), we model

\begin{equation}
  V_{ijp}
  \;=\;
  \beta_0
  + \boldsymbol{\beta}_1^\top \mathbf{d}^{\textsf{spec}}_i
  + \boldsymbol{\beta}_2^\top \mathbf{d}^{\textsf{test}}_i
  + \boldsymbol{\beta}_3^\top \mathbf{d}^{\textsf{code}}_i
  + \boldsymbol{\beta}_4^\top \mathbf{d}^{\textsf{path}}_p
  + \gamma_j
  + \varepsilon_{ijp},
  \label{eq:mixedlogit}
\end{equation}

where $\gamma_j \sim \mathcal{N}(0, \sigma_{\textsf{sample}}^2)$ is a random intercept for sample and $\varepsilon_{ijp}$ is residual error. Reference category qwen3.6 was chosen as an interior mono baseline on the $V$ scale (M3 sits in the middle of the six mono cells on strict-AND valid-closure fraction; note M3 is the strongest mono on the $\bar{Y}$ scale, but the reference selection is by $V$-scale interiority since $V$ is the outcome). The full Bernoulli GLMM (logit link on $\Pr[V_{ijp} = 1]$) did not converge on our data volume ($n \approx 5{,}000$ rows across 150 samples with unbalanced strata); we report the linear-probability approximation of Eq.~\ref{eq:mixedlogit} with mixed-effects standard errors, following the standard fallback when GLMM convergence fails. Coefficient vectors $\hat{\boldsymbol{\beta}}_1, \hat{\boldsymbol{\beta}}_2, \hat{\boldsymbol{\beta}}_3$ quantify each stage's marginal contribution to valid-closure probability relative to placing qwen3.6 there; $\hat{\boldsymbol{\beta}}_4$ quantifies path main effects.

\subsubsection{Judge--human agreement, sensitivity, capability preservation}
\label{sec:experiment_setup_supporting_methods}

Judge--human agreement (Cohen's $\kappa$, Krippendorff's ordinal $\alpha$),
threshold-sensitivity ranges $\tau \in \{0, 0.5, 0.8\}$ and $\rho \in \{2, 3, 4\}$
with Kendall $\tau_{\textsf{rank}}$-based stability check
(Section~\ref{sec:experiment_setup_sensitivity}), the post-hoc decision-reason
and filter-reason labels applied to every closure row, and the strongest-mono-
baseline restriction underlying the capability-preservation rate --- these
supporting protocols are formally defined in
Appendix~\ref{sec:appendix_methodology} to keep this section focused on the
core statistical instruments. Results depending on them are reported in
Section~\ref{sec:results}. The 60-pair human-evaluation protocol
(Section~\ref{sec:experiment_setup_human_eval}) uses the same 0--4 rubric
as the judge SLM.

\subsubsection{Sensitivity analysis}
\label{sec:experiment_setup_sensitivity}

To defend against threshold-choice attacks, closure rates are reported at
multiple $(\tau, \rho)$ combinations (Path~1 $\tau \in \{0, 0.5, 0.8\}$;
judge $\rho \in \{2, 3, 4\}$). Ordinal stability of the per-cell ranking
across thresholds is quantified using Kendall's $\tau_{\textsf{rank}}$
between the reference $(\tau = 0, \rho = 3)$ and each alternative. Full grid
in Appendix~\ref{sec:appendix_threshold}.

\subsection{Human-evaluation protocol (summary)}
\label{sec:experiment_setup_human_eval}

To validate the framework's judge signal against human raters, we
pre-registered a 60-pair five-annotator study drawn from the completed
sweep. Sampling is stratified into three 20-pair buckets (frontier,
agreement, disputed) around the metric--judge agreement structure.
Annotators rate each pair on the judge SLM's $0$--$4$ rubric (4~=~identical,
3~=~equivalent, 2~=~approximately equivalent, 1~=~clearly different,
0~=~unrelated), blinded to cell identity, closure path, judge rating, and
bucket assignment. Analysis computes Krippendorff's ordinal $\alpha$ across
annotators, pairwise linear- and quadratic-weighted Cohen's $\kappa$,
Cohen's $\kappa$ between the judge SLM and the majority human rating,
within-1 agreement rates, and an FPR/FNR ablation of the three candidate
decision policies. Full protocol in Appendix~\ref{sec:appendix_human_eval_protocol};
results in Section~\ref{sec:results_human_eval}.


